\Needspace{5\baselineskip}
\SourceLeaf{429}%
\begin{center}10.\end{center}
\noindent\hspace*{1.2em}Before further transforming the expression just found, we shall make a few remarks about it.

A normal to the curve in its plane corresponds to two directions on the sphere, depending on which side of the curve it is drawn. The direction towards which the curve is \textls[100]{concave} is denoted by $\lambda'$, and the other by the opposite point on the sphere. These two points, like $L$ and $\mathfrak{L}$, are separated from $\lambda$ by $90^{0}$ and thus lie on a great circle. Since $\mathfrak{L}$ is also separated from $\lambda$ by $90^{0}$, we have $\mathfrak{L}L=90^{0}-L\lambda'$ or $=L\lambda'-90^{0}$, and therefore
\[
\cos L\lambda'=\pm\sin\mathfrak{L}L,
\]
where $\sin\mathfrak{L}L$ is necessarily positive. Since in our analysis $r$ is taken to be positive, the sign of $\cos L\lambda'$ will be the same as that of
\[
Z(\xi\xi T+2\xi\eta U+\eta\eta V)
\]
and a positive value of the latter quantity therefore indicates that $L\lambda'$ is less than $90^{0}$, or that the curve is concave on the side on which the projection of the normal to the curved surface onto the plane lies. A negative value indicates that the curve is convex on that side. We may therefore also set, in general,
\[
\frac1r=\frac{Z(\xi\xi T+2\xi\eta U+\eta\eta V)}{\sin\mathfrak{L}L}
\]
if we take the radius of curvature to be positive in the first case and negative in the other. $\mathfrak{L}L$ is now the angle that our cutting plane makes with the plane tangent to the curved surface. We see that in different cutting planes through the same point and the same tangent, the radii of curvature are proportional to the sine of the inclination. Because of this simple relation, we shall henceforth restrict ourselves to the case in which this angle is a right angle and the cutting plane therefore passes through the normal to the curved\space%
